\section{Recovery Experiment and Certificate Protocol}\label{sec:r44-protocol}
\paragraph{Scaling inputs.}
For $k$ branches, set $b_j=1/5+3j/[5(k-1)]$ for $j=0,\ldots,k-1$. Weights are proportional to $1+(j\bmod5)$, curvatures are $(1+(j\bmod9))/3$, and $r=2$, $q=1$. The promise is $4\bar b/5$. Catalogs are uniform with denominator $N-1$. Charges are $\rho(c)=ac$ with the slope $a$ displayed in Table~\ref{tab:r44-work}; zero charges have no hidden fixed fee. Cases zero through eight use price tolerance $10^{-3}$. Cases nine through eleven hold the 32-branch, 17-level, two-symbol, zero-risk instance fixed and vary that tolerance through $10^{-1}$, $10^{-3}$, and $10^{-5}$. All twelve declared configurations are retained.

\paragraph{Measurement.}
The run records wall time with a monotone performance clock and peak Python-traced allocation with \texttt{tracemalloc}. These times include instrumentation. Peak traced allocation is not total process resident memory. The maximum numerator or denominator bit length stored in price dynamic-programming tables is recorded; it is not a claim to measure every temporary integer in the interpreter. The separate checker is timed after certificate serialization. The JSON environment record includes Python and operating-system versions; local and continuous-integration executions are not pooled as repeated statistical samples.

\paragraph{Independent checks.}
Each certificate stores model and catalog input, opening charges, both prices, complete endpoint policies, mixture weight, fixed union policy, a separately feasible original-budget policy, original upper bound, price error, and charge excess. The checker recomputes price upper bounds with explicit branch-support maximization in the independent R43 checker, not the recovery solver's edge-construction module. It verifies that the endpoint payoff plus its priced promise residual equals the respective global price upper bound. The union support, branch targets, pre-draw tiers, probabilities, each realization, root sum, payoff, and every installed charge are checked directly as fractions. Finally it checks the union identity, target mixture, price error, charge excess, and the distinct original-budget interval. A malformed record fails rather than being accepted with a numerical tolerance.

\paragraph{Test coverage and limits.}
The 144 exhaustive comparisons use the fixed seed 2026092444. The corruption suite changes, in turn, a reported upper bound, endpoint bound, endpoint payoff, price, mixing weight, branch target, intermediate tier, lottery normalization, installed level, paid charge, charge excess, price error, budget, and original lower bound. All must be rejected. Five invalid-input cases test nonpositive accuracy, floating-point input, zero budget, negative tolerance, and an infeasible promise. These tests complement the proof; they do not establish a same-budget approximation theorem. Nor does the observed union size establish a universal bound smaller than $2m$.

\paragraph{Reproduction.}
From the repository root, run \nolinkurl{revisions/or-r44-resource-augmentation-20260924/code/tests.py}, \nolinkurl{revisions/or-r44-resource-augmentation-20260924/code/study.py}, and \nolinkurl{revisions/or-r44-resource-augmentation-20260924/code/inherited.py} with Python. The checker accepts the JSON files in this revision's \texttt{results/certificates} directory. Table generation reads only the executed JSON outputs. Source, evidence, reader hashes, and the preservation comparison are recorded separately so that a successful build is not confused with mathematical globality.
